
\documentclass[12pt]{article}
\usepackage{amsmath}
\usepackage{amsfonts}
\usepackage{amssymb}
\usepackage{geometry}
\geometry{a4paper, margin=0.8in}
\pagenumbering{gobble}

\title{Homework assignment for 02YMECH \\ Chapter: Kinematics}
\begin{document}
\date{}
\maketitle
\vspace{-0.8in}
\begin{enumerate}
    \item A particle is under uniform circular motion (counter-clockwise) with a radius of 5 m and a constant speed of 10 m/s. The center of the circle is at point (3m, 4m) in the xy-plane. 
    \begin{enumerate}
      \item Write the position vector of the particle as a function of time, assuming it starts at the point (8m, 4m) at time t = 0.
      \item Calculate the velocity and acceleration vectors at $t=\pi/2$ seconds and show that they are perpendicular to each other by considering their dot product. Is this true for all time?
    \end{enumerate}


    \item The position vectors of two particles are given by:
    \[\mathbf{r}_1(t) = (1 + 2t)\,\hat{\imath} + (3 - t)\,\hat{\jmath}\]
    \[\mathbf{r}_2(t) = (4 - t)\,\hat{\imath} + (1 + 3t)\,\hat{\jmath}\]
    \begin{enumerate}
        \item Find the velocity and acceleration of each particle.
        \item Find the time at which the two particles are closest to each other and determine that minimum distance.
    \end{enumerate}
    \item A stationary body of total mass $M = 13\ \text{kg}$ suddenly disintegrates into three fragments with masses
\[
m_1 = 3\ \text{kg}, \qquad m_2 = 4\ \text{kg}, \qquad m_3 = 6\ \text{kg}.
\]
Right after the break-up:
\begin{itemize}
    \item fragment $m_1$ moves due east with speed $v_1 = 4\ \text{m/s}$,
    \item fragment $m_2$ moves due north with speed $v_2 = 3\ \text{m/s}$.
\end{itemize}
Using only conservation of linear momentum, determine the velocity (both speed and direction) of fragment $m_3$ immediately after the break-up. 

\item A block of mass $m = 4~\mathrm{kg}$ rests on a horizontal surface. The coefficient of static friction between the block and the surface is $\mu_s = 0.5$, and the coefficient of kinetic friction is $\mu_k = 0.3$.

A horizontal force $F(t)$ is applied to the block such that it increases linearly with time:
\[
F(t) = 2t \quad \text{(in N, with $t$ in seconds)}.
\]

\begin{enumerate}
    \item Determine the time at which the block starts to move.
    \item Once the block starts moving, find its acceleration immediately after it starts moving.
    \item Find the velocity of the block 5 seconds after it starts moving. 

\textit{(Assume $g = 10~\mathrm{m/s^2}$.)}
\end{enumerate}
          
\end{enumerate}
\end{document}
